\chapter{Multi-tip Langmuir Probe}

\section{What's a Langmuir probe}

\begin{figure}
    \caption{Diagram of simple langmuir probe.}
\end{figure}

\section{Hardware}

\subsection{Shaft Components}

\subsubsection{Garages}

\begin{figure}
    \caption{Figure 2.6 from Joe's thesis}
\end{figure}

\subsubsection{Shaft materials}

\subsubsection{Centering Collet}

\subsubsection{Feed-thru}

\subsection{Tip Components}

\begin{figure}
    \caption{Picture of end of probe}
\end{figure}

\subsubsection{Probe body}

\begin{figure}
    \caption{Probe body from engineering sketch}
\end{figure}

\subsubsection{Tip material + size}

\paragraph{Why use Mo}

\paragraph{Why use the size we chose}

\subsubsection{Tip Shielding}

\paragraph{Why shield at all?}

\paragraph{First order mach effects to avoid}

\subsection{Bdot Coils}

\subsubsection{How many wraps}

\subsubsection{Approximate area}

\subsection{Length Scales}

\begin{table}
    \caption{Length scales of probe}
\end{table}

\section{Circuitry}

\subsection{Bias Source}

\subsubsection{Manually Controlled}

\paragraph{Principles of operation}

- Equation to use for bias calculation

\begin{figure}
    \caption{Circuit diagram of old setup. Add switches.}
\end{figure}

\begin{figure}
    \caption{Image of old setup}
\end{figure}

\paragraph{Problems}

- Had to be manually controlled for drastic bias changes
- Wires often contacted each other or parts would lose contact

\subsubsection{Remote Controlled}

\paragraph{Principles of operation}

- Equation to use for calculation

\begin{figure}
    \caption{Circuit diagram of new setup. Add switches.}
\end{figure}

\begin{figure}
    \caption{Image of new setup}
\end{figure}

\begin{table}
    \caption{Bias circuit component values}
\end{table}

\paragraph{Problems}

- Destroyed a power supply
- A resistor was designed and thus burnt out

\paragraph{Calibration}

\subsection{Probe Circuit}

\subsubsection{Principles of operation}

\begin{figure}
    \caption{Simple circuit diagram}
\end{figure}

\paragraph{Plasma draws current}

\paragraph{Voltage drop across sense resistor}

\paragraph{Measure high frequency voltage drop}

\subsubsection{Single-ended measurement}

\begin{figure}
    \caption{Image of old circuit board}
\end{figure}

\subsubsection{Differential measurement}

\begin{figure}
    \caption{Image of new circuit board}
\end{figure}

\begin{table}
    \caption{Probe circuit component values}
\end{table}

\subsubsection{Circuit Solution}

\paragraph{How do we solve this}

\subparagraph{Kirkoff's laws}

\subparagraph{System of equations}

\subparagraph{SymPy}

\paragraph{equations}

\paragraph{Assumptions}

\subsubsection{Calibration}

\paragraph{Calibrating in isat}

\paragraph{Calibrating on the exponential}

\paragraph{Bayesian calibration}

\paragraph{Calibration using pressure wave}

\subsection{Bdot coils}

\subsubsection{In-shaft circuit}

\paragraph{Why do we have that?}

\paragraph{Where is it placed approximately?}

\subsubsection{Circuit model}

\subsubsection{Calibration}

\paragraph{Directional}

\paragraph{Transfer Function}

\section{Plasma Theory}

\subsection{Simple Plasma Theory (Hutchinson)}

\subsubsection{Derivation}

\subsubsection{Pros \& Cons}

\subsection{BRL}

\subsubsection{Idea behind it}

\subsubsection{Approximations of it}

\paragraph{Planar probe}

\paragraph{Chen Approximation}

\begin{table}
    \caption{Chen parameters}
\end{table}

\paragraph{Even more approximation}

\subsection{Effect of electric fields}

\subsubsection{Where does the electric field come from}

\paragraph{Majority inductive}

\paragraph{Some static}

\subparagraph{What is the expected static magnitude}

\subsubsection{What does it do to the tip measurements}

\paragraph{Theoretical understanding}

\begin{figure}
    \caption{What does an electric field do to the tip measurements diagram.}
\end{figure}

\paragraph{Physical reality}

\begin{figure}
    \caption{Example plot where E splitting tips}
\end{figure}

\subsection{Minimum measurable density}

\subsubsection{What's the debye Length}

\subsubsection{Sheath scale}

\subsubsection{Tip-tip sheath overlap}

\section{Curve Fitting}

\subsection{Least squares Fitting}

\subsubsection{Benefits (easy)}

\subsubsection{Issues (not representative of error)}

\subsection{ODR}

\subsubsection{What does error look like in our signals}

\subsubsection{How does ODR take that into account}

\subsubsection{How to apply parameter bounds}

\paragraph{Our method}

\paragraph{Other options}

\subsection{Monte-carlo Fitting}

\begin{figure}
    \caption{Diagram of local minimums}
\end{figure}

\subsubsection{What kind of fits do we get if we just trust a local minimum}

\subsubsection{How do we sample initial conditions}

\subsubsection{More robust methods for getting global minimums}

\subsection{Combining multiple shots}

\subsubsection{Using some time points from each shot}

\subsubsection{Combining all IV points and fitting that}

\subsubsection{Problems}

\paragraph{Timing alignment}

We'll go into this more in the PA section.

\paragraph{Variability of plasma parameters}

\section{Results}

What is the benefit of doing all this work?

\begin{figure}
    \caption{Example fit using many shots, multiple local parameter fits, and errors}
\end{figure}

\subsection{Measurement of low densities}

\subsubsection{How does Hutchinson vs BRL treat ion saturation}

\subsubsection{What are the experimental predictions that Hutchinson vs BRL have}

\begin{figure}
    \caption{Plot of isat currents from BRL}
\end{figure}

\subsubsection{How does that change the results}

\begin{figure}
    \caption{Radial chord of initial density with Hutchinson vs BRL.}
\end{figure}

\subsection{Noise resilience}

Not sure how best I should show this.

\subsection{Fits more accurate}

- Each fit uses many data points

- Don't have to arbitrarily cutoff data points

\subsection{Extracting electric fields}

\subsubsection{Where do we expect electric fields to be}

\subsubsection{Do we see them there}

\begin{figure}
    \caption{Comparison of expected E from Ohm's law and measured E.}
\end{figure}
